%! Departures from Linearity — Unit 2, Lesson 2.5 — Warm-Up
% ONE PAGE, blank and key. Three items; do not add a fourth.
% GRADUAL RELEASE: the warm-up activates prior knowledge before the guided notes.
%   Item 1 — spiral to Lessons 2.3–2.4: compute and read a residual (the notes lean on residuals
%            all period).
%   Item 2 — two residual plots, one random and one curved: which says "a line is right"? Seeds
%            the Residual Plots row, where r = -0.94 sits beside a U-shaped plot (the crux).
%   Item 3 — where would one new point move the line most? Seeds leverage and influence in the
%            Unusual Points row. Leave the class's answer on the board.
% Every display is pre-drawn; nothing is sketched.
\documentclass[12pt]{article}
\usepackage{apstats-article}
\usepackage{apstats-boxes}

\begin{document}

\pageheader{Unit 2, Lesson 2.5}{Warm-Up}
% NAMESTRIP: no \namedateperiod here — the name row lives on the cover only.

\vspace{0.28in}

\begin{enumerate}[leftmargin=1.9em, itemsep=0.42in]

  \item \textbf{From last lesson.} A city's taxi fares follow the line
        $\hat{y} = 2.50 + 0.80x$, where $x$ is the miles driven and $\hat{y}$ is the predicted
        fare in dollars. A 10-mile ride cost \$12.00.

        \vspace{12pt}
        Predicted fare: \blank{2.4cm} \qquad Residual: \blank{2.4cm}
        \vspace{16pt}

        Did the line predict too high or too low for this ride? \blank{2.8cm}

  \item Two residual plots from two different lines. Which one says a line is a good model
        for its data? Explain.

        \vspace{10pt}
        \begin{center}
        \begin{tikzpicture}[scale=0.62]
          \draw[linegray] (0,0) -- (7,0);
          \draw[->, thick] (0,-1.8) -- (0,2);
          \node[left, font=\scriptsize] at (0,0) {0};
          \foreach \x/\y in {0.5/0.6,1.1/-0.9,1.6/1.2,2.2/-0.3,2.8/0.8,3.3/-1.3,3.9/0.2,
                             4.4/1.0,5.0/-0.6,5.5/0.4,6.1/-1.1,6.6/0.9}
            \fill[navy] (\x,\y) circle (0.09);
          \node[font=\bfseries\small] at (3.5,2.4) {Plot A};
          \node[font=\scriptsize] at (7.3,0) {$x$};
          \begin{scope}[xshift=9.5cm]
            \draw[linegray] (0,0) -- (7,0);
            \draw[->, thick] (0,-1.8) -- (0,2);
            \node[left, font=\scriptsize] at (0,0) {0};
            \foreach \x/\y in {0.5/1.6,1.1/0.8,1.6/0.1,2.2/-0.5,2.8/-0.9,3.3/-1.1,3.9/-1.0,
                               4.4/-0.7,5.0/-0.2,5.5/0.4,6.1/1.0,6.6/1.7}
              \fill[navy] (\x,\y) circle (0.09);
            \node[font=\bfseries\small] at (3.5,2.4) {Plot B};
            \node[font=\scriptsize] at (7.3,0) {$x$};
          \end{scope}
        \end{tikzpicture}
        \end{center}
        \vspace{6pt}
        \writeline\\[18pt]\writeline

  \item The line below was fit to the six navy points. \textbf{One} more point will be added:
        either P or Q. Which one would change the \emph{slope} of the line more? Why?

        \vspace{10pt}
        \begin{center}
        \begin{tikzpicture}[xscale=0.62, yscale=0.42]
          \draw[->, thick] (0,0) -- (13.5,0) node[below, font=\scriptsize] {$x$};
          \draw[->, thick] (0,0) -- (0,5.5) node[left, font=\scriptsize] {$y$};
          \draw[navy, thick] (0.5,1.0) -- (6.5,4.0);
          \foreach \x/\y in {1/1.4,2/1.6,3/2.4,4/2.7,5/3.5,6/3.6}
            \fill[navy] (\x,\y) circle (0.1 and 0.15);
          \fill[redacc] (3.5,4.9) circle (0.13 and 0.19) node[right, font=\bfseries\small] {\,P};
          \fill[redacc] (12.5,1.2) circle (0.13 and 0.19) node[right, font=\bfseries\small] {\,Q};
        \end{tikzpicture}
        \end{center}
        \vspace{6pt}
        \writeline\\[18pt]\writeline

\end{enumerate}

\end{document}
